\uuid{i3Py}
\exo7id{4499}
\titre{exo7 4499}
\auteur{quercia}
\organisation{exo7}
\datecreate{2010-03-14}
\isIndication{false}
\isCorrection{true}
\chapitre{Série numérique}
\sousChapitre{Familles sommables}
\cpge{mpsi-mp2i.procedes-sommatoires-discrets.a, pcsi.series-numeriques.a, ptsi.series-numeriques.a, mpsi-mp2i.procedes-sommatoires-discrets.d, mp-mpi.series-numeriques-vectorielles.a, pt.series-numeriques.b, mp-mpi.series-numeriques-vectorielles.b, pc.series-numeriques, psi.series-numeriques.a, pt.series-numeriques.a}
\module{Analyse}
\niveau{L1}
\difficulte{}

\contenu{
\question{
On réordonne les termes de la série harmonique alternée en prenant tour
à tour $p$ termes positifs puis $q$ termes négatifs, $p,q\ge 1$.
Calculer la somme de la série correspondante.
}
\reponse{
$\ln2 + \frac12\ln(p/q)$.
}
}
